\chapter{Glossary of Key Terms}

\noindent This glossary covers Pali terms that appear frequently in the endnotes. Terms are listed in Roman alphabetical order of the Pali. Page references point to the first significant discussion.

\bigskip

\begin{description}[style=nextline, labelwidth=3.5cm, leftmargin=3.8cm, font=\itshape]

\item[abhidhamma] ``Higher teaching.'' The analytical and classificatory literature of the Pali Canon; the third ``basket'' of the Tipiṭaka. The Netti's method is closely related to Abhidhamma analysis.

\item[adhiṭṭhāna] ``Resolution, determination.'' As a mode (\textit{hāra}), it distinguishes general categories (``oneness'') from specific distinctions (``difference'') within a single term.

\item[anattā] ``Not-self.'' One of the three characteristics of all conditioned phenomena, along with impermanence and suffering.

\item[anicca] ``Impermanent.'' The first of the three characteristics.

\item[anusaya] ``Underlying tendency.'' A latent defilement that has not yet risen to the level of active obsession but lies dormant, ready to be triggered.

\item[āvaṭṭa] ``Turning back, conversion.'' As a mode, it converts passages between the defilement-side and the purification-side.

\item[avijjā] ``Ignorance.'' The root of the view-character's defilements; the failure to see the Four Noble Truths.

\item[bhāvanā] ``Development, cultivation.'' The practice of training the mind, particularly through meditation.

\item[bhavaṅga] ``Factor of existence, life-continuum.'' The constituents that sustain continued existence.

\item[bodhipakkhiyā dhammā] ``Qualities conducive to awakening.'' The thirty-seven factors grouped into seven sets: four foundations of mindfulness, four right strivings, four bases of power, five faculties, five powers, seven awakening-factors, and the Noble Eightfold Path.

\item[catubyūha] ``Fourfold array.'' As a mode, it determines the actual meaning intended by the speaker through four analytical questions.

\item[cetanā] ``Intention, volition.'' The mental factor that drives action (\textit{kamma}).

\item[citta] ``Mind, consciousness, thought.'' Used broadly for the knowing aspect of experience.

\item[dassana] ``Seeing, vision.'' In the Netti's terminology, the initial breakthrough of insight at stream-entry.

\item[desanā] ``Teaching, instruction.'' As a mode, it identifies the gratification-danger-escape structure in any passage.

\item[dhamma] ``Teaching, truth, mental quality, phenomenon.'' Context determines which meaning applies. As a foundation of mindfulness, it refers to mental qualities or principles.

\item[diṭṭhi] ``View.'' Can be wrong view (\textit{micchādiṭṭhi}) or right view (\textit{sammādiṭṭhi}).

\item[diṭṭhicarita] ``View-character.'' A person dominated by ignorance, prone to wrong views, self-mortification, and annihilationism or eternalism.

\item[dukkha] ``Suffering, unsatisfactoriness, pain.'' The first Noble Truth.

\item[hāra] ``Mode, method of reading.'' Literally ``carrying away, removing.'' The Netti's sixteen modes of textual analysis.

\item[hārasampāta] ``Modes Applied, convergence of modes.'' The section where all sixteen modes are demonstrated on a single passage.

\item[hāravibhaṅga] ``Mode Analysis.'' The detailed exposition of each of the sixteen modes.

\item[hetu] ``Cause.'' Distinguished from \textit{paccaya} (condition) by being specific and unshared.

\item[indriya] ``Faculty, controlling principle.'' The five spiritual faculties: faith, energy, mindfulness, concentration, wisdom.

\item[jhāna] ``Absorption, meditative attainment.'' The four stages of deep concentration.

\item[khandha] ``Aggregate, group.'' The five aggregates of clinging: form, feeling, perception, formations, consciousness. Also: the three training-aggregates (virtue, concentration, wisdom).

\item[kilesa] ``Defilement, mental corruption.'' The unwholesome mental factors that obstruct liberation.

\item[lakkhaṇa] ``Characteristic, mark.'' As a mode, it identifies the defining characteristic of each term.

\item[magga] ``Path.'' The Noble Eightfold Path; the fourth Noble Truth.

\item[mūlapada] ``Root term.'' The eighteen master categories of the Netti's classification system.

\item[naya] ``Method, strategy, approach.'' The five methods of the Netti: Delight-Round, Lion's Play, Direction-Gazing, Goad, and Triple-Turning.

\item[neyya] ``What is to be guided, one who needs to be led.'' One of three types of learner.

\item[nibbāna] ``Extinguishment, liberation.'' The cessation of suffering; the unconditioned.

\item[niddesavāra] ``Description section.'' The third of the Netti's four main divisions.

\item[nirodha] ``Cessation.'' The third Noble Truth; the cessation of craving.

\item[niyyāna] ``Leading out.'' The characteristic shared by all awakening-qualities.

\item[otaraṇa] ``Entry, descent.'' As a mode, it enters a passage through five analytical doorways: faculties, dependent origination, aggregates, elements, sense-bases.

\item[paccaya] ``Condition.'' Distinguished from \textit{hetu} (cause) by being shared and generic.

\item[padaṭṭhāna] ``Basis, standing-place, proximate cause.'' As a mode, it traces each mental state to its supporting condition.

\item[paññatti] ``Designation, concept.'' As a mode, it identifies over twenty types of designation (teaching-, origin-, comprehension-, abandonment-designation, etc.).

\item[paññā] ``Wisdom, discernment.'' The faculty of understanding things as they are.

\item[parikkhāra] ``Requisite, equipment.'' As a mode, it maps the chain of causes and conditions.

\item[pariyuṭṭhāna] ``Active obsession.'' A defilement that has risen from its latent state (\textit{anusaya}) to active expression.

\item[paṭiccasamuppāda] ``Dependent origination.'' The twelve-link chain of conditions that sustains the cycle of existence.

\item[paṭiniddesavāra] ``Detailed exposition section.'' The fourth and largest of the Netti's four main divisions.

\item[sakkāyadiṭṭhi] ``Identity-view.'' The view that there is a self among the five aggregates. Twenty-fold when applied to all five aggregates in four modes.

\item[samādhi] ``Concentration, unification of mind.'' One of the three training-aggregates.

\item[samāropana] ``Integration.'' As a mode, it weaves together the results of all other modes into a unified reading.

\item[samatha] ``Calm, tranquility.'' One of the two wings of meditation, paired with \textit{vipassanā} (insight).

\item[saṃkilesa] ``Defilement, corruption.'' The dark side of the teaching; threefold: craving-defilement, view-defilement, misconduct-defilement.

\item[saṃsāra] ``The round of rebirth, continued existence.'' The cycle of birth, death, and re-becoming.

\item[saṅkhāra] ``Formations, volitional constructions.'' The fourth aggregate; also the second link of dependent origination.

\item[sāsanapaṭṭhāna] ``Foundation of the Teaching.'' The final major section of the Netti, classifying suttas into sixteen types.

\item[sīla] ``Virtue, moral conduct.'' The first of the three training-aggregates.

\item[sodhana] ``Clearing.'' As a mode, it distinguishes genuine answers from restatements of the question.

\item[sutta] ``Discourse, teaching.'' A text attributed to the Buddha or a senior disciple.

\item[taṇhā] ``Craving, thirst.'' The origin of suffering; the second Noble Truth. Threefold: sense-craving, existence-craving, non-existence-craving.

\item[taṇhācarita] ``Craving-character.'' A person dominated by desire, prone to sensual indulgence and eternalism.

\item[ugghaṭitaññū] ``One who understands on the spot.'' The keenest type of learner.

\item[upādāna] ``Clinging, fuel.'' Fourfold: sense-pleasure clinging, existence-clinging, view-clinging, self-doctrine clinging.

\item[vevacana] ``Synonym.'' As a mode, it identifies when different terms in different suttas point to the same reality.

\item[vibhatti] ``Analysis, sorting.'' As a mode, it separates mixed passages into their wholesome and unwholesome components.

\item[vicaya] ``Investigation, analysis.'' As a mode, it investigates the terms, structures, and components of a passage.

\item[vijjā] ``True knowledge.'' The opposite of \textit{avijjā} (ignorance).

\item[vimokkha / vimokkhamukha] ``Liberation / liberation-door.'' Three doors: emptiness (\textit{suññatā}), signless (\textit{animitta}), desireless (\textit{appaṇihita}).

\item[vipassanā] ``Insight.'' The practice of seeing things as they really are; one of the two wings of meditation.

\item[vipallāsa] ``Distortion, perversion.'' Fourfold: seeing permanence, pleasure, self, and beauty where none exist. Operates at three levels: perception, thought, and view.

\item[vodāna] ``Purification, cleansing.'' The bright counterpart of defilement.

\item[yuttihāra] ``Fitness Mode.'' Tests whether an interpretation is consistent with the suttas, the discipline, and dependent origination.

\end{description}
